%\documentclass[gray,handout, pdftex, 11pt]{beamer}
%\documentclass[handout, pdftex, 11pt]{beamer}
\documentclass[pdftex, 11pt]{beamer}

\usepackage{pgfpages}
\usepackage[utf8]{inputenc}
\usepackage[T1]{fontenc}
\usepackage{lmodern}
\usepackage[italian]{babel}
\usepackage{graphicx}
\usepackage{microtype}
\usepackage{acronym}
\usepackage{array}
%\usepackage{natbib}
\usepackage{tikz-uml}

\usepackage{tikz}
\usetikzlibrary{intersections, arrows, shapes, decorations.pathreplacing, decorations.pathmorphing, calc}

\usepackage{appendixnumberbeamer}

\mode<presentation>{
  %-------------------------1
  \usetheme{Boadilla}
  \usecolortheme{beaver}
  %-------------------------1
  %-------------------------2
  %\usetheme{Goettingen}
  %\usecolortheme{sidebartab}
  %-------------------------2
  %\useoutertheme[right]{sidebar}
  %\usefonttheme{default}
  \setbeamercovered{transparent}
  %\setbeameroption{show notes on second screen=right}
  \setbeamertemplate{navigation symbols}{}

  \bibliographystyle{abbrv}  
  %\renewcommand\bibfont{\scriptsize}
  \setbeamertemplate{bibliography item}{\textbullet}
  \setbeamertemplate{itemize item}{\checkmark}
  \setbeamertemplate{itemize subitem}{-}
  \setbeamertemplate{enumerate items}[default]
  \setbeamertemplate{sections/subsections in toc}[square]
}

%stili
\tikzstyle{stato}=[fill=yellow, circle, draw=black!80, minimum size=2cm]
\tikzstyle{ftato}=[fill=yellow!50, circle, draw=black!80, dashed, minimum size=2cm]
\tikzstyle{freccia}=[->, very thick, >=stealth', draw=black!80]

\title[Progress report]{\textbf{PROGRESS REPORT}}
\subtitle{tkLayout developers meeting}
\institute[CERN]{
  %{\Large\textbf{CERN}}\\{European Organization for Nuclear Research}\\[0.5cm]
  %\\[0.2cm]
  European Organization for Nuclear Research\\[0.5cm]
  \includegraphics[width=2cm]{img/LogoBadge.pdf}\\
}

\author[Stefano Martina]{
  %\\[0.2cm]
  \textbf{Stefano MARTINA}\\
  {\small stefano.martina@cern.ch}
}

\date[\today]{\flushright \today}

\begin{document}

\begin{frame}[plain,noframenumbering]
  \titlepage
\end{frame}

\begin{frame}
  \begin{itemize}
  \item \alert{Novit\`a} nelle mappe di irradiazione
    \pause
    \begin{itemize}
    \item Cambiata \alert{struttura} delle righe del file
      (funzionalit\`a gi\`a gestita)
      \pause
    \item Presenza di un \alert{header}
      \pause
    \item Presenza di \alert{pi\`u mappe} con risoluzioni diverse
    \end{itemize}
    \pause
  \item I valori vengono impostati \alert{leggendoli} dall'header
    \pause
  \item Vi \`e un \alert{gestore} che \alert{contiene} le mappe e che
    si occupa di usare quella con risoluzione migliore per ogni modulo
  \end{itemize}
\end{frame}

\begin{frame}
  \begin{center}
    \begin{tikzpicture}
      \uncover<2>{
        \umlclass[x=0,y=0]
        {IrradiationMapsManager}
        {- irradiationMaps : vector<IrradiationMap>}
        {
          + addIrradiationMap(mapFile : string)\\
          + getIrradiationPower(coord : pair<double,double>) : double
        }
      }
      \umlclass[x=0,y=-4.5]
      {IrradiationMap}
      {
        - irradiation : vector<vector<double>{}>\\
        - rMin : double\\
        - rMax : double\\
        $\cdots$
      }
      {
        + ingest(mapFile : string)\\
        + operator < (confronted : IrradiationMap) : bool\\
        + getIrradiation(coord : pair<double,double>) : double\\
        + isInRegion(coord : pair<double,double>) : bool\\
        $\cdots$
      }
      \uncover<2>{
        \umluniassoc[arg2=irradiationMaps,mult2=*]{IrradiationMapsManager}{IrradiationMap}
      }
    \end{tikzpicture}
  \end{center}
\end{frame}
\end{document}
